\section{Conclusion}
\label{sec:conclusion}

We presented Multi-Scale Embodied Memory (\Acronym{}), an approach for equipping VLAs with long-horizon memory. \Acronym{} enables VLAs to perform long-horizon tasks that require tens of minutes of memory, while obeying real-world latency constraints. \Acronym{} also enables in-context adaptation of manipulation strategies and, when combined with the \pizerosix{} VLA, achieves state-of-the-art performance across a wide range of manipulation tasks. To achieve this, \Acronym{} introduces a mixed-modal memory architecture that combines short-horizon, video-based memory, with a long-horizon, language-based memory mechanism. 
We believe that \Acronym{} is only the first step towards building robot policies that can effectively manage very long-horizon memory. Future work can explore how we can scale memory to last beyond the horizon of a single episode, to span weeks, months, or years of deployment, and allow us to build robots that learn continually at deployment time.

